\chapter[INTRODUÇÃO]{INTRODUÇÃO}

Plantas de processo contínuo não são conjuntos de malhas independentes. Um reator, um separador e uma coluna de \textit{stripping} interconectados por reciclos formam um sistema acoplado em que a ação de qualquer controlador local propaga efeitos ao longo de toda a cadeia de processo — às vezes amplificados, às vezes com atraso suficiente para enganar o operador. Nesse contexto, a pergunta central não é "como sintonizar este PID?", mas sim "quem decide como a planta deve operar?" — e essa pergunta não tem resposta no nível de controle regulatório. Ela exige uma camada supervisória com visão holística da planta \cite{Larsson2000, Skogestad2004}. O \textbf{Processo Tennessee Eastman} (TEP), proposto por Downs e Vogel em 1993, é o \textit{benchmark} canônico que cristaliza exatamente esse problema \cite{DOWNS1993}: cinco unidades interconectadas, reações não-lineares acopladas, laço de reciclo e 20 distúrbios programados — uma planta deliberadamente impossível de controlar sem coordenação supervisória, e amplamente adotada na literatura precisamente por isso.

Larsson e Skogestad, revisando décadas de literatura sobre controle plant-wide, apontam que a maioria das teorias de controle disponíveis pressupõe uma estrutura de controle já dada e, por isso, falha em responder às perguntas fundamentais que Foss formulou em 1973: quais variáveis controlar, quais medir, quais manipular e como conectá-las \cite{Larsson2000}. Controle plant-wide significa definir a \textit{filosofia de controle} da planta como um todo, com ênfase nas decisões estruturais: seleção das variáveis manipuladas, seleção das medições, decomposição do problema em subproblemas e configuração das conexões entre eles. O tipo de controlador — PID, MPC, LQG — é apenas o último item dessa lista \cite{Skogestad2004}. Skogestad formaliza esse raciocínio em um procedimento sistemático com duas direções complementares: \textit{top-down}, partindo dos objetivos operacionais e dos graus de liberdade disponíveis para atendê-los; e \textit{bottom-up}, construindo a partir da camada de estabilização regulatória. A hierarquia resultante separa escalas de tempo distintas — controle regulatório (segundos), supervisório (minutos) e otimização em tempo real (horas) — cada camada usando os setpoints da camada inferior como seus graus de liberdade \cite{Larsson2000, Skogestad2004}. Aplicado ao TEP, esse procedimento revela que a planta possui 8 graus de liberdade em regime estacionário: no modo nominal, 5 deles ficam consumidos por restrições economicamente ativas, restando 3 graus livres cuja escolha de variáveis controladas determina a perda econômica total quando distúrbios chegam \cite{Larsson2000}. 

O controle plant-wide, como Larsson e Skogestad o definem é um problema de \textbf{integração}. Coordenar uma camada regulatória (segundos), uma supervisória (minutos) e uma de otimização (horas) sobre uma planta com reciclos e acoplamentos exige que componentes heterogêneos troquem informação, negociem setpoints e mantenham consistência de estado mesmo diante de falhas e distúrbios. Este trabalho é desenvolvido sob a ênfase de \textbf{Integração de Sistemas de Controle}, com o objetivo de implementar essa hierarquia de controle sobre o TEP usando técnicas modernas de integração de sistemas.

A resposta formal a esse problema é o conceito de \textbf{Sistema Cyber-Físico} (CPS): a integração entre computação e processos físicos, onde computadores e redes monitoram e controlam o processo por meio de laços de realimentação em que o físico afeta a computação e vice-versa \cite{Hehenberger2016}. Hehenberger et al.\ apontam que uma abordagem de engenharia baseada em CPS impõe dois passos metodológicos: modelar explicitamente cada subsistema e integrar esses modelos em um sistema global coerente \cite{Hehenberger2016}. A \textit{parte cyber} de um CPS é precisamente a rede de integração — a camada que interconecta subsistemas, torna seus estados visíveis e permite que comandos fluam de volta ao processo. Esse requisito de integração horizontal e vertical é o que torna necessária uma camada de comunicação industrial com semântica padronizada e, consequentemente, motiva as escolhas tecnológicas deste trabalho.

Para a camada supervisória desse CPS — o componente que observa o estado da planta e age para manter a operação dentro de parâmetros declarados — a literatura cloud-native oferece um candidato bem estudado. Shuiguang et al.\ \cite{Shuiguang2023} caracterizam computação cloud-native como a combinação de microserviços, containerização e orquestração automática de componentes distribuídos, com Kubernetes como o orquestrador open-source que se tornou padrão de fato na indústria. O mecanismo central que o torna atraente para esse papel é o \textit{reconciliation controller loop}: uma comparação contínua entre o estado desejado (\textit{Spec}) e o estado observado (\textit{Status}), seguida de ações para convergir um ao outro \cite{Burns2016}. Essa estrutura — declarar política, observar realidade, reconciliar continuamente — guarda uma semelhança estrutural não trivial com o papel que Skogestad atribui à camada supervisória: manter variáveis controladas em setpoints ótimos usando os graus de liberdade da camada regulatória como alavancas. Este trabalho explora a hipótese de que esse padrão arquitetural pode servir como orquestrador supervisório de um CPS industrial baseado no TEP.

O presente trabalho constrói três artefatos como ponto de partida para esse universo de possibilidades. O primeiro é um microserviço da planta TEP implementado em Rust — uma refatoração do modelo original em FORTRAN para uma linguagem de alta performance, exposto via gRPC. O segundo é uma Interface Homem-Máquina (IHM) web em FastAPI com WebSocket, funcionando como cliente dessa planta e permitindo monitoramento em tempo real e injeção dos distúrbios. O terceiro é um Operator Kubernetes básico que observa a planta, demonstrando a viabilidade de usar o padrão declarativo do Kubernetes como orquestrador supervisório sobre componentes industriais. A interface entre os três componentes é mapeada segundo o padrão OPC~UA, respeitando a norma IEC~62541 \cite{IEC62541_1}. O objetivo não é implementar estratégias de controle nem estabelecer lógica supervisória definitiva — é produzir esses artefatos iniciais como fundação técnica e ponto de partida concreto para os trabalhos futuros que o referencial teórico delineia.

\section{PROBLEMA}
\label{sec:problema}

O problema abordado neste trabalho é a implementação de controle plant-wide sobre o Processo Tennessee Eastman como um desafio de integração industrial. A TEP, com suas unidades interconectadas e laços de reciclo, exige uma camada supervisória que coordene componentes heterogêneos em diferentes escalas de tempo. Para que esse sistema seja extensível e interoperável — e não apenas um protótipo isolado —, seus componentes devem expor interfaces segundo normas industriais estabelecidas, de modo que possam ser compostos, observados e substituídos por ferramentas compatíveis. A abordagem adotada é a de um Sistema Cyber-Físico com Kubernetes como orquestrador supervisório candidato, o que impõe que tanto a planta simulada quanto os demais componentes respeitem uma semântica de integração padronizada.

\section{A PROPOSTA}
\label{sec:proposta}

Este trabalho propõe refatorar o modelo original do TEP de FORTRAN para um microserviço em Rust, modelado segundo a hierarquia de integração da norma IEC~62264 \cite{IEC62264} e com variáveis de processo expostas conforme o padrão OPC~UA \cite{IEC62541_1}. Sobre essa base, são implementados os seguintes componentes:

\begin{enumerate}
  \item Simulação da planta TEP em Rust com integração numérica RK4, comunicação via gRPC e modelagem compatível com a hierarquia Purdue definida pela IEC~62264 \cite{IEC62264};
  \item Interface Homem-Máquina (IHM) web aderente à norma IEC~63303 \cite{IEC63303}, funcionando como cliente da planta para monitoramento em tempo real e injeção de distúrbios;
  \item Implementação dos distúrbios programados descritos no artigo original de Downs e Vogel \cite{DOWNS1993};
  \item Verificação da fidelidade dinâmica da simulação por meio da análise do comportamento da planta sob os distúrbios implementados;
  \item Implementação de um \textit{Custom Resource Definition} (CRD) Kubernetes conectando um contêiner Kind à planta, validando a arquitetura básica do CPS proposto.
\end{enumerate}

\section{OBJETIVOS}
\label{sec:objetivos}

\subsection{Objetivo Geral}

Implementar uma simulação funcional do Processo Tennessee Eastman em Rust — refatoração do modelo original em FORTRAN — e integrá-la em uma arquitetura de Sistema Cyber-Físico com componentes heterogêneos interoperáveis via protocolos abertos, servindo como base técnica para o desenvolvimento futuro de estratégias de controle plant-wide.

\subsection{Objetivos Específicos}
\begin{enumerate}
    \item Verificar que a simulação em Rust reproduz fielmente a dinâmica da planta TEP descrita por Downs e Vogel, por meio da análise do comportamento sob os distúrbios programados do artigo original.
    \item Verificar que a IHM opera corretamente como cliente da planta, respeitando a interface OPC~UA e os requisitos de monitoramento em tempo real definidos pela IEC~63303.
    \item Demonstrar que os componentes heterogêneos do sistema — simulação, IHM e orquestrador — são interoperáveis via a interface padronizada, sem acoplamento direto entre implementações.
    \item Validar que o padrão CRD do Kubernetes consegue observar o estado da planta, estabelecendo a viabilidade arquitetural do Kubernetes como orquestrador supervisório de um CPS industrial.
    \item Estabelecer uma base técnica documentada e reprodutível que permita trabalhos futuros evoluir para estratégias completas de controle plant-wide sobre a mesma infraestrutura.
\end{enumerate}

\section{ORGANIZAÇÃO DO TRABALHO}

Este documento está organizado da seguinte forma. O Capítulo~2 apresenta os fundamentos teóricos: o Processo Tennessee Eastman, os princípios de controle plant-wide, o diagnóstico de qualidade de malhas, a analogia cloud-native, o IEC~61499, o padrão de Operators do Kubernetes, o protocolo gRPC e o OPC UA. O Capítulo~3 descreve o desenvolvimento da infraestrutura proposta, detalhando cada componente e as decisões de arquitetura. O Capítulo~4 apresenta os resultados dos 17 experimentos realizados e a análise do comportamento da planta sob distúrbios. O Capítulo~5 conclui o trabalho e apresenta as direções de trabalhos futuros, incluindo a incorporação de métricas de qualidade de malha à camada supervisória.
